\section{Identification}

In this section we describe the phases that led to the acquisition of the domain
knowledge, the identification of the requirements to be implemented, and the
definition of the reasoning strategy.


\subsection{Knowledge Acquisition}

The acquisition of the knowledge of the application domain was a complex and
laborious task. It was carried out by drawing on three different kinds of source:

\begin{itemize}
    \item \emph{Books and manuals}:
    \begin{itemize}
        \item \emph{DSM-5-TR. Diagnostic and Statistical Manual of Mental
        Disorders}~\cite{dsm}.
        \item \emph{The ABC of Psychopathology. Exploration, identification and
        treatment of mental disorders}~\cite{abc}.
    \end{itemize}

    \item \emph{Web resources}:
    \begin{itemize}
        \item \url{http://www.ipsico.org/default.htm}
        \item \url{http://www.cpsico.com/index.htm}
        \item \url{http://www.ansia-depressione.net}
    \end{itemize}

    \item \emph{Human expert}: the consultancy provided by a psychologist.
\end{itemize}

From the study of the diagnostic criteria we learned that each ``parent'' anxiety
disorder can be specialised into several ``child'' anxiety disorders.

Although the DSM-5-TR is a sound diagnostic manual, it does not allow a
knowledge engineer to encode exactly the various disorders that fall
under each category. Indeed, while it reports the diagnostic criteria related to
each ``parent'' disorder together with the possible specialisations, it does not
contain all the information needed to differentiate the diagnoses among the
``child'' disorders.

To address this issue, in addition to the use of manuals and psychology portals,
we relied on a human expert who---thanks to their experience in terms of
knowledge, heuristics and reasoning mechanisms---made it possible to simplify the
complexity of the work.

Starting from the diagnostic criteria of each ``parent'' anxiety disorder, we
carried out a process of refinement and specialisation of the various criteria so
as to identify every single ``child'' anxiety disorder. At the end of the
knowledge-acquisition work, 43 ``child'' anxiety disorders were identified; they
are reported in Table~\ref{tab:disorders}.

\begin{longtable}{p{0.34\linewidth} p{0.58\linewidth}}
\caption{Classification of the anxiety disorders handled by the expert
system.}\label{tab:disorders}\\
\toprule
\textbf{Category} & \textbf{Disorders} \\
\midrule
\endfirsthead
\toprule
\textbf{Category} & \textbf{Disorders} \\
\midrule
\endhead
\midrule
\multicolumn{2}{r}{\footnotesize\itshape continued on next page}\\
\endfoot
\bottomrule
\endlastfoot

Panic Disorder &
Panic Disorder Without Agoraphobia; Panic Disorder With Agoraphobia. \\
\addlinespace

Agoraphobia Without History of Panic Disorder &
Agoraphobia Without History of Panic Disorder. \\
\addlinespace

Specific Phobia &
Specific Phobia, ``Animal'' Type; Specific Phobia, ``Natural Environment'' Type;
Specific Phobia, ``Blood-Injection-Injury'' Type; Specific Phobia,
``Situational'' Type; Specific Phobia, ``Other'' Type. \\
\addlinespace

Social Phobia &
Generalised Social Phobia; Circumscribed Social Phobia. \\
\addlinespace

Obsessive-Compulsive Disorder &
Obsessive-Compulsive Disorder with Aggressive Content (and with Poor Insight);
with Contamination Content (and with Poor Insight); with Sexual Content (and with
Poor Insight); Hoarding Type (and with Poor Insight); with Religious Content (and
with Poor Insight); Order and Symmetry Type (and with Poor Insight); Somatic Type
(and with Poor Insight); Superstitious Type (and with Poor Insight); Perfection
and Responsibility Type (and with Poor Insight); with Miscellaneous Content (and
with Poor Insight). \\
\addlinespace

Post-Traumatic Stress Disorder &
Acute Post-Traumatic Stress Disorder; Acute Post-Traumatic Stress Disorder with
Delayed Onset; Chronic Post-Traumatic Stress Disorder; Chronic Post-Traumatic
Stress Disorder with Delayed Onset. \\
\addlinespace

Acute Stress Disorder &
Acute Stress Disorder. \\
\addlinespace

Generalised Anxiety Disorder &
Generalised Anxiety Disorder. \\
\addlinespace

Anxiety Disorder Due to a General Medical Condition &
With Generalised Anxiety; With Panic Attacks; With Obsessive-Compulsive
Symptoms. \\
\addlinespace

Substance-Induced Anxiety Disorder &
With Generalised Anxiety; With Panic Attacks; With Obsessive-Compulsive Symptoms;
With Phobic Symptoms. \\
\end{longtable}


\subsection{Requirements}

In the following sections we provide a brief description of the requirements that
the system must satisfy and of the features to be implemented.


\subsubsection{Programming Language}

The expert system must be implemented in Prolog.

Prolog, an acronym for \emph{PROgramming in LOGic}, is a very powerful and
flexible logic programming language. It is particularly well suited to Artificial
Intelligence and to non-numerical programming.

By using Prolog as the programming language, it is possible to express the
knowledge about the identified objects and relations in the form of facts and
rules. It is also possible to ask questions of the user so as to derive further
information useful to the inference process.

Two behavioural features required by the system are not provided natively by
Prolog: the handling of uncertainty and an explanation module. They are discussed
below.


\subsubsection{Handling of Uncertainty}

The first behavioural feature we want to implement is the handling of
uncertainty.

Prolog reasons in a categorical way, in the sense that every element can only be
either true or false, with no other possibility.

In the medical field, this kind of reasoning turns out to be too restrictive. The
diagnosis of disorders or diseases may, in fact, rest on rather uncertain
knowledge, both because symptoms and causes may be only probably correlated with
a disorder or disease, and because the patient may not answer with complete
certainty.

The handling of the uncertainty intrinsic to the available knowledge becomes even
more necessary in a critical domain such as the diagnosis of anxiety disorders,
where the mental disorder must be managed as quickly and correctly as possible, in
order to prevent the exacerbation of the symptoms from ``destroying'' the overall
functioning of the patient or from significantly interfering with their
occupational or social functioning.

From these premises, the need emerges to move away from the categorical nature of
assertions and to converge towards a probabilistic nature, in which it is possible
to assign a degree of certainty to every known fact or rule.

This notion of certainty is not to be understood in the mathematical sense of the
term, but rather as the subjective confidence one has in the truth of an assertion
or in the fact that an inference rule turns out to be applicable.

Indeed, it is still a matter of debate whether it is more appropriate to use
probability theory in its most rigorous sense or in a more informal and
simplified version.

In the project under examination, we chose to use the more informal and
simplified version for two simple reasons. First, because we want the expert
system to emulate human reasoning, which is certainly intuitive and not very
mathematical. Second, because one cannot disregard the greater computational
simplicity of an intuitive and less rigorous method.


\subsubsection{Explanation Module}

The second behavioural feature concerns the implementation of an explanation
module.

In an expert system, the explanation capability is of fundamental importance.
Users, in fact, tend not to trust an artificial system slavishly, nor to accept
blindly the solution of a problem produced in an uncontrollable way.

This capability is even more necessary in uncertain domains, such as that of
medical diagnosis, in order to increase the trust that the physician or the expert
places in the advice of the system, and to enable them to discover a possible
error in the reasoning process.

In keeping with this requirement, the expert system under examination provides
explanations about its own way of proceeding on two occasions: when it is asked
\emph{why} a given question is being posed, and when it is asked \emph{how} the
solution of a problem was reached.


\subsection{Reasoning Strategy}

As far as the type of inference to be used is concerned, the choice fell on
backward chaining.

Backward chaining is a type of inference in which one starts from a possible
solution of the problem and tries to verify, by applying the production rules in
reverse, whether there exists a rule capable of providing that solution. It is
opposed to forward chaining, in which one starts from the known facts contained in
the knowledge base and applies the production rules forward in order to deduce new
facts. Like the handling of uncertainty, backward chaining is not provided
natively by Prolog and must be implemented by the inference engine.

The inference process will moreover use the \emph{left-most} computation rule and
the \emph{depth-first} search strategy. Thus, in the inference process, given a
query

\[
\texttt{:-}\ G_1,\, G_2,\, \dots,\, G_n
\]

the left-most literal will always be selected. Given the literal to be resolved
(for example, $G_1$), the first clause whose head is unifiable with $G_1$ will
always be selected.

As for the behaviour that the inference process must adopt upon the occurrence of
possible negative answers, we opted for a ``non-failure'' policy. Considering, in
fact, that the expert system has to handle the uncertainty of the information, we
decided to process negative answers as positive answers, but with a
``complemented'' degree of certainty. For instance, a negative answer with a
degree of certainty of $0.7$ for a feature required as true will be processed by
the system as a positive answer with a degree of certainty of $0.3$. In this way,
it is possible to bring every computation to completion, returning the degree of
certainty with which the conclusion turns out to be true.

Finally, since some diagnostic criteria require the verification of undesirable
features, the system must provide the possibility of specifying, for a given goal,
which sub-goals are required as true and which as false. Depending on the case, it
will be the inference engine that correctly computes the degree of certainty.

The reasoning strategy introduced in this section will be taken up again and
documented in depth in the sections devoted to the description of the inference
engine.
